Forecasting is a critical capability for global decision-making, spanning domains from energy grid management to geopolitical strategy. Traditionally, human forecasters have been relied upon for qualitative, event-based reasoning, while statistical models have dominated numerical time-series analysis. However, human cognitive bandwidth is severely limited in high-volume data tasks, and human experts are often subject to heuristic biases in low-data regimes. In this paper, we evaluate the forecasting performance of Large Language Models (LLMs) across these two extremes of the forecasting spectrum. Using \methodname, we conduct a unified assessment on short-term high-data numerical forecasting (\ettm) and long-term low-data event-based forecasting (\forecastbench). Our results demonstrate that LLMs possess strong zero-shot capabilities in both domains. In high-data tasks, LLMs process 96-hour historical windows with error magnitudes (MSE 3.988) comparable to naive statistical baselines, far exceeding human manual capacity. In low-data event forecasting, \methodname achieves a Brier score of 0.246, outperforming human expert ``superforecasters'' (0.287) and matching the aggregated human crowd (0.243). These findings establish LLMs as versatile and robust forecasters capable of matching or surpassing human expertise across disparate data regimes.
